% !TeX program = pdflatex
% ALIUS Bulletin native reconstruction source.
% Compile from the ALIUS-bulletin project root. This file does not include a pre-existing article PDF.
\providecommand{\ALIUSRootPrefix}{}

\ifdefined\ALIUSIssueBuild
\else
\documentclass[a4paper,11pt]{article}
\usepackage[margin=22mm]{geometry}
\usepackage[T1]{fontenc}
\usepackage[utf8]{inputenc}
\usepackage[hidelinks]{hyperref}
\usepackage[protrusion=true,expansion=false]{microtype}
\usepackage{parskip}
\fi
\title{Body meets self}
\author{Frederique de Vignemont interviewed by Raphael Milliere \& Carlota Serrahima}
\date{ALIUS Bulletin 4}

\ifdefined\ALIUSIssueBuild
\else
\begin{document}
\fi
\maketitle
\section*{Metadata}
\textbf{Piece type:} interview\par
\textbf{DOI:} 10.34700/gza7-yy49\par
\textbf{Keywords:} bodily awareness, sense of body ownership, pain, touch, peripersonal space, somatoparaphrenia, depersonalization, deafferentation, rubber hand illusion, peripersonal space\par
\section*{Abstract}
In this interview, Frederique de Vignemont discusses her wide-ranging and influential research program on philosophical issues related to bodily awareness. The conversation explores core questions of this research program, such as the existence of a sense of body ownership, the nature of pain and touch, and the role of the peripersonal space, as well as methodological questions regarding the role of empirical evidence in philosophical investigation and the value of arguments from phenomenal contrast in philosophy of mind. In the course of this discussion, Vignemont defends her interpretation of pathological conditions such as somatoparaphrenia, depersonalization disorder, and peripheral deafferentation, as well as experimental bodily illusions such as the rubber hand illusion.\par
\section*{Extracted Text}
\begingroup
\small
\noindent ALIUS Bulletin n4 (2020) aliusresearch.org/bulletin Body Meets Self Abstract In this interview, Frederique de Vignemont discusses her wide-ranging and influential research program on philosophical issues related to bodily awareness. The conversation explores core questions of this research program, such as the existence of a sense of body ownership, the nature of pain and touch, and the role of the peripersonal space, as well as methodological questions regarding the role of empirical evidence in philosophical investigation and the value of arguments from phenomenal contrast in philosophy of mind. In the course of this discussion, Vignemont defends her interpretation of pathological conditions such as somatoparaphrenia, depersonalization disorder, and peripheral deafferentation, as well as experimental bodily illusions such as the rubber hand illusion. keywords: bodily awareness, sense of body ownership, pain, touch, peripersonal space, somatoparaphrenia, depersonalization, deafferentation, rubber hand illusion, peripersonal space Your research on bodily awareness is one of the most exhaustive and influential in the field of empirically informed philosophy of mind. Bodily awareness is also the topic to which you have devoted the most significant part of your academic production, which covers a very wide range of philosophical issues related to it - from the specificity of particular bodily experiences, such as pain or touch (Vignemont, 2017b), to our capacity for empathy for the bodily sensations of others (Vignemont, 2017a). Within this range of issues, one that is central in your work, and in particular in your recent book Mind the Body (Vignemont, 2018a), is the status of bodily awareness as a form of self-consciousness. The book is indeed a deep exploration of the so- called sense of bodily ownership, namely of how subjects experience their body An interview with Frederique de Vignemont by Raphael Milliere \& Carlota Serrahima Frederique de Vignemont frederique.de.vignemont@ens.fr Ecole Normale Superieure, Paris, France Raphael Milliere rm3799@columbia.edu Columbia University, New York, USA Carlota Serrahima carlota.serrahima@uantwerpen.be University of Antwerp, Antwerp, Belgium Cite as: de Vignemont, F., Milliere, R. \& Serrahima, C. (2020). Body Meet Self. An interview with Frederique de Vignemont by Raphael Milliere and Carlota Serrahima. ALIUS Bulletin, 4, 107-133 https://doi.org/10.34700/gza7-yy49 Vignemont - Body Meet Self ALIUS Bulletin n4 (2020) aliusresearch.org/bulletin 108 as their own. What motivated you to work on bodily awareness, and on bodily self-consciousness in particular, in the first place? I wrote my PhD thesis on self-consciousness, and more specifically on immunity to error through misidentification. At the time, bodily self- awareness was just one among other forms of self-awareness and the sense of agency was a lot more popular in the philosophical and the experimental community than the sense of body ownership. Which made it also more interesting. While many were discussing newly found results on action attribution in schizophrenia and delusions of control, there was simply nothing in philosophy on disorders of bodily self-awareness such as somatoparaphrenia (the delusion that one's limb doesn't belong to oneself) or on the rubber hand illusion (the illusion that a rubber hand is one's own hand). Even in cognitive science, these two minimal forms of self-awareness were left disconnected. When I went to work with Professor Patrick Haggard at University College London, my objective was to investigate how agency and ownership interacted. My project, however, ended up working only on the bodily side. The fact is that Queen's Square was a unique place for such a topic, with everyday bringing of new fascinating results. I was lucky to witness and actively participate in the emergence of a new field of research, with a new "body" community that, in fifteen years, was to become so important. For every new exploratory research, there was still very little theory and much conceptual confusion - the perfect challenge for a philosopher... Philosophers' concerns about self-consciousness have classically covered further aspects of our consciousness of ourselves. One of these aspects is, paradigmatically, whether and how phenomenal awareness of psychological states amounts to self-consciousness proper. For instance, the famous Cartesian dictum states the alleged evidence of the existence of (one's)self on the grounds of states we would now call cognitive, even when in doubt about the existence of the body. Analytic philosophy has of course produced an amazing body of literature addressing this psychological dimension of self- consciousness. Investigations on bodily awareness in contemporary philosophy of mind, however, sometimes run parallel to the insights in this body of literature. In what sense, if any, do you think these two topics, and correspondingly these two areas of research, complement each other? Vignemont - Body Meet Self ALIUS Bulletin n4 (2020) aliusresearch.org/bulletin 109 The body is a material entity located in space and time in the same way as a rock, a tree or a bird. Yet we do not perceive and experience our body like those other objects. What makes it unique is that it bears a special relation to the self, and to self-awareness. However, I want to make it clear that: (i) all bodily awareness does not have to be self-awareness and (ii) not all forms of self-awareness have to be embodied. In my work, I have repeatedly emphasized the fact that one can experience bodily sensations in a part of one's body though one does not experience this body part as being one's own. This is the case in some patients suffering from somatoparaphrenia or from depersonalization. But also, this is most probably the case of some animals that do not have self-referential abilities even at the nonconceptual level. Put it another way, bodily awareness is more widespread in the animal kingdom than self-awareness. My second claim targets embodied theories of the self, which I do not endorse. The point is not to deny that bodily awareness may be a gateway to self-awareness at the developmental level. But the fact that the body may be at the origin of self-awareness does not entail that every single instance of self-awareness is constitutively embodied. One might reply that in sensory deprivation self-awareness becomes quite thin. This, however, does not suffice to show that self-awareness consists in bodily awareness. What might be necessary for self-awareness is not specifically information about the body, but incoming information in general. The case of sensory deprivation, however, cannot dissociate the two factors. Furthermore, discussions over the notion of immunity to error through misidentification relative to the first- person reveal a clear difference between psychological and bodily self- awareness. Since Evans (Evans 1982), most agree that the self-ascription of bodily judgments is immune to error if grounded in the right way of gaining information about one's body. However, there seems to be a difference between the judgment "I think" and the judgment "my legs are crossed". Immunity directly follows from the peculiarities of mental states but it does not follow from the nature of bodily states. One should not reduce the former to the latter. Vignemont - Body Meet Self ALIUS Bulletin n4 (2020) aliusresearch.org/bulletin 110 One of the main distinctive features of your approach is its substantive reliance on empirical data, both from experimental and clinical psychology. Cases such as the rubber hand illusion, somatoparaphrenia, or xenomelia (a disorder characterized by a desire to be disabled or having discomfort with being able-bodied), are now well-known in philosophical discussion of bodily awareness, and this is surely in part because of your own effort to highlight their theoretical import. On the grounds of some of these cases you draw conclusions, for instance, on the ontology of consciousness: on your view, postulating the existence of certain bodily feelings is the best way to explain some features of the empirical cases, against a general principle of phenomenal parsimony (Vignemont, 2018a, section 1.1, 2019a). Your work on the epistemology of bodily awareness appeals to empirical findings as well: they help evaluate the validity of the thesis of the immunity to error through misidentification of bodily self-ascriptions (Vignemont, 2011, 2018a, chap. 3). What would you say is, in general, the methodological value of empirical cases for philosophers? And what would you say is the value of a priori philosophical arguments nowadays, in particular for research on bodily awareness? In my work, I constantly go back and forth between the philosophical and the empirical literature. Cognitive science has grown to such an extent that to neglect what it has to tell us about the mind seems a pity. However, this is not to say that experimental findings are always relevant to philosophical issues. Nor is it to say that one can simply endorse whatever conclusions cognitive neuroscientists reach. The first step is to decide whether the question that one has in mind can be addressed empirically. For instance, my current research concerns the nature of valence, whether it should be explained in terms of content or in terms of attitude. No experimental result can answer this question. By contrast, I have also been working on spatial awareness. There, the discovery of the specific multisensory and motor properties of peripersonal space in cognitive The fact that the body may be at the origin of self-awareness does not entail that every single instance of self-awareness is constitutively embodied. " " Vignemont - Body Meet Self ALIUS Bulletin n4 (2020) aliusresearch.org/bulletin 111 neuroscience is crucial. This discovery shows that we do not process our immediate surroundings in the same way as far space. The second step for philosophers is to proceed to a task of conceptual clarification: to what extent does peripersonal space differ from egocentric space, from reaching space, or from personal space? Do all studies and all paradigms investigate the same notion? (Vignemont et al., forthcoming) etc. Only once the notion is clearer can one draw the philosophical implications of the findings. For instance, there has been much discussion on enactivism and sensorimotor theories of perception. What new light do these results shed on these theories? How is perception related to action in close space (Vignemont, 2018c)? What is important to keep in mind in philosophy of cognitive science is that at no point should a hypothesis rely exclusively on empirical data. We have recently seen enough that one cannot always replicate results and one should always be very cautious when using them, especially when one is a philosopher with little statistical expertise. Experimental findings can be the starting point but they do not replace philosophical arguments. Ideally, they should be the cherry on top, which only empirically confirm what has been found conceptually. And the other way around: in your experience of collaboration with neuroscientists and psychologists (e.g., Folegatti et al., 2012; Gouzien et al., 2017), what do you think is the specific contribution that philosophers can make to the cognitive sciences? I have had the chance to collaborate with many wonderful researchers from various fields, including cognitive psychologists, like Alessandro Farne and Patrick Haggard, neuroscientists like Tamar Makin and Tania Singer, psychiatrists, like Victor Pitron and Adrienne Gouzien, and even roboticists like Aldo Faisal and Silvestro Micera. Some of our work together was theoretical. For instance, with Tamar Makin and Silvestro Micera, we just Experimental findings can be the starting point but they do not replace philosophical arguments. Ideally, they should be the cherry on top, which only empirically confirm what has been found " " Vignemont - Body Meet Self ALIUS Bulletin n4 (2020) aliusresearch.org/bulletin 112 have a short opinion piece accepted in TICS on augmentative technology (Makin et al., 2020). We have combined our respective expertise to determine what the most promising path is for the integration of prostheses, arguing for what we call soft embodiment, which we define at the conceptual, neural and computational levels. Other contributions have been more directly experimental. Ideally, some philosophical theories should at least be empirically testable and can be a source of inspiration for cognitive scientists. Furthermore, by constantly anticipating objections, philosophers are well prepared to think of control conditions. Beside direct collaboration, I have found it very interesting to discuss work in progress with junior researchers. As a philosopher, I have more time for reading than experimentalists and this may be helpful to bring the results back into the big picture. To defend the claim that there is a phenomenology of bodily ownership, you rely in large part on so-called arguments from phenomenal contrast. This label was coined by Siegel (2007), although the relevant class of arguments has a long history in philosophy. Arguments from phenomenal contrast are generally used to arbitrate introspective disagreements regarding the existence of a specific kind of phenomenology, by proceeding in two steps: 1 It is argued that a pair of similar experiences E1 and E2 differ - in some small but noticeable way - with respect to their phenomenal character (i.e., what it is like to have them). 2 It is argued that the best explanation of the phenomenal contrast between E1 and E2 is that one experience involves a phenomenal feature F that the other lacks. In recent years, arguments from phenomenal contrast have been notably used in the debate on the existence of a sui generis cognitive phenomenology (e.g., Strawson, 1994; Chudnoff, 2015), and the debate over whether 'high-level' properties - such as the property of being a pine tree - are represented in visual experience (Bayne, 2009; Siegel, 2010). In your own work, you use arguments from phenomenal contrasts in two complementary ways (Vignemont, 2007, 2013, 2018a, forthcoming). Firstly, you argue that the phenomenal contrast between the experiences produced by the rubber hand illusion experiment in the synchronous (illusory) and asynchronous (non-illusory) conditions, respectively, is best explained by the following hypothesis: subjects feel a phenomenology of ownership over the rubber hand in the synchronous condition, while they lack such Vignemont - Body Meet Self ALIUS Bulletin n4 (2020) aliusresearch.org/bulletin 113 phenomenology in the asynchronous condition. If this hypothesis is true, then the phenomenology of bodily ownership exists at least in rare experimental conditions. Secondly, you argue that the phenomenal contrast between a healthy individual's experience of her limb and a patient's experience of her limb in what you call 'disownership syndromes' - which include most prominently somatoparaphrenia, but also depersonalization disorder and other conditions - is best explained by the following hypothesis: in the former (ordinary) condition, but not in the latter (pathological) conditions, one's experience of one's limb involves a phenomenology of bodily ownership. If this hypothesis is true, then the phenomenology of bodily ownership exists not only in rare experimental conditions, but is also prevalent in the ordinary experience of healthy individuals. Arguments from phenomenal contrast face a well-known challenge: if there is a genuine phenomenal contrast between E1 and E2, it might really be the case that E1 lacks a phenomenal feature F that E2 has; but it could also be the case that E1 has an additional feature F' that E2 lacks. With respect to the phenomenology of bodily ownership, it could be the case that (a) the experience of the rubber hand in the asynchronous condition of the rubber hand illusion, and (b) the experience of the affected limb in disownership syndromes, do not lack a feature that the contrasted condition has, but instead involve an additional and abnormal sense of alienation with respect to the relevant limb (Billon \& Kriegel, 2015; Chadha, 2018). In fact, there is no shortage of reports from somatoparaphrenic and depersonalized patients that describe a limb as 'alien' or 'strange', or use a lexicon that seems neutral with respect to the two interpretations of the phenomenal contrast (e.g., by saying that the limb is 'dead' or 'rotten'). Do you think this challenge can be addressed, such that the debate over the interpretation of the phenomenal contrast can be settled with a reasonable degree of confidence? If so, what do you think is the most promising kind of evidence that could rule out the alternative interpretation? First, I should point out that the method of phenomenal contrast is far from perfect. It has come upon heavy criticisms for classic cases (as in the visual experience of pine trees), and it is even more controversial for bodily awareness. As discussed in my latest paper on the phenomenal contrast of ownership, the "all things being equal" rule cannot apply here when it comes to the sense of bodily ownership because the cases we have, such as the rubber hand illusion and disownership syndromes, involve many other differences besides ownership (Vignemont, 2020). That's why, since my Analysis paper in Vignemont - Body Meet Self ALIUS Bulletin n4 (2020) aliusresearch.org/bulletin 114 2013, I have always said that a single case taken in isolation is not sufficient to show that there is a phenomenology of ownership and thus to reject the deflationist view (Vignemont, 2013). What we need is to consider a range of cases taken together and abstract what they have in common. But even then, at most what we can offer is an inference to the best explanation. This means that it definitely cannot rule out alternative explanations. What I argue is just that the hypothesis of a phenomenology of ownership can easily account for a range of cases. Now the hypothesis that by default there is no phenomenology of ownership (hereafter default hypothesis) is a valid alternative interpretation of the evidence, but is it a better one? To settle the debate, we need to know what the arguments are in favour of this view. One could argue that it fits more our introspective reports. As repeatedly emphasized by philosophers from all sides, bodily ownership is not phenomenologically salient under normal condition. The simplest explanation would be that it is because there is no phenomenology of ownership. However, bodily awareness is recessive in general, this is not a problem specific to ownership. Even more generally, what is too familiar always goes to the background of consciousness. And what is more familiar than the fact that this is our own body? There is a real question here but it is an issue about the rules of consciousness in general, and it is not specific to the sense of ownership. Another argument in favour of the default hypothesis is that it seems more parsimonious. But is it really? After all, it involves that in the non-default cases, bodily experiences can represent non-ownership (Chadha, 2018). Advocates of a conservative conception of perception already complain that ownership is a too high-level property for being part of the bodily content, but non-ownership fares even worse. The fact that it is only in some rare cases does not make it more admissible. As mentioned above, your main argument from phenomenal contrast in favor of the existence of a phenomenology of bodily ownership in the ordinary experience of healthy individuals relies on empirical evidence regarding somatoparaphrenia (SP) (Vignemont, 2007, 2013, 2018a, forthcoming). SP is a monothematic delusion (typically caused by a brain lesion) characterized by the patients' belief that one of their body parts is not really Vignemont - Body Meet Self ALIUS Bulletin n4 (2020) aliusresearch.org/bulletin 115 theirs. A number of patients also believe that the affected limb belongs to someone else. Your argument from SP starts with the assumption that there is a phenomenal contrast between (a) what it is like for an SP patient to have bodily sensations in the body part whose ownership they deny (e.g., their right hand), and (b) what it is like for a healthy individual to have bodily sensations in the corresponding body part (e.g., their right hand). SP is associated with a large number of severe motor and somatosensory impairments, which frequently include unilateral neglect, hemiplegia on the contralesional side of the body (paralysis of half of the body, including the affected limb), impairment of the ability to determine the position of one's affected limb through proprioception, hemianaesthesia on the contralesional side (impairment or loss of tactile perception), and hemianopia on the contralesional side (loss of vision in half of the visual field) (Vallar \& Ronchi, 2009; Romano \& Maravita, 2019). In many cases, the ability of SP patients to feel bodily sensations at all in their affected limb is heavily impaired. Arguments from phenomenal contrast normally focus on 'minimal pairs' of experiences that are as similar as possible, with the exception of one clear and unique difference. Thus, one might expect the argument from SP to focus on the phenomenal contrast between two experiences involving the same type of bodily sensation (e.g., touch) in the same limb (e.g., the subject's right hand), to bring out a single phenomenal feature that is missing in one experience and present in the other, all else being equal. However, the difference between the bodily experiences of somatoparaphrenic patients and that of healthy individuals are so dramatic that it might be difficult to find such a minimal pair of experiences. In Mind the Body, you acknowledge that the ability to feel touch is frequently affected in SP (2018a, p. 40), but you point out that many SP patients can feel pain in the affected limb, and that in a handful of rare cases these patients can also feel and report touch in the affected limb (e.g., Bottini et al., 2002). You argue that the argument from SP can be set up by using these kinds of cases to bring out the relevant phenomenal contrast. However, there is little evidence that in such cases, the determinate experience of pain (or touch) of an SP patient and the determinate experience of pain (or touch) of a healthy individual differ, if at all, with respect to their phenomenal character. While there is ample evidence that overall, SP patients have a range of abnormalities in the way they can experience their affected limb, there is less evidence to support the specific claim that patients with intact ability to feel pain or touch Vignemont - Body Meet Self ALIUS Bulletin n4 (2020) aliusresearch.org/bulletin 116 in the affected limb have a different determinate phenomenology from healthy individuals when they experience pain or touch in the affected limb. Do you agree with this assessment of available empirical evidence? If so, to what extent do you think the first step of your argument from somatoparaphrenia - establishing the existence of a phenomenal contrast between a minimal pair of bodily experiences - might be affected at all by this assessment? This is a problem encountered by all attempts to use the phenomenal contrast method: can the contrast be explained by other things? Now in the specific case of ownership, I have collected descriptions not only of somatoparaphrenia, but also of peripheral deafferentation, which is especially interesting, I believe, because it is as pure of a comparison as it can be. Deafferented patients have no brain lesion. This avoids the risk of reasoning deficit, neglect, and other attentional or cognitive disorders, which might be found in brain-lesioned patients or psychiatrist patients. The patients suffered only from a peripheral loss of proprioception and touch. But they can still feel pain and thermal sensations. Now one of them, Ian Waterman, describes how at the beginning, he did not feel his body as being his own (Cole 1995). This, however, did not last. We thus have a contrast between the beginning of the disease and a later stage. One way to describe it is to say that at the early stage, Ian feels pain in some legs and that at the late stage, he feels pain in his legs. The comparison is not between a patient and healthy subjects, but intra-individual. Now the difference between the two stages is not only a matter of ownership. There is also an agentive contrast. Because of the loss of proprioception, deafferented patients need to learn to exploit vision to replace proprioception to control their bodily movements. So, at the early stage, they have no control over their body and at the late stage, they have regained it. I do not think that the agentive contrast shows that there is no difference at the level of the preserved bodily sensations. Instead, I believe that it can explain the ownership contrast: this temporary loss of agency impacts their preserved bodily sensations, thus explaining the loss of ownership. To conclude, as I have said again and again, I do not believe that somatoparaphrenia suffices to answer all our questions. What we need to do Vignemont - Body Meet Self ALIUS Bulletin n4 (2020) aliusresearch.org/bulletin 117 is to consider all the cases that are relevant. And it is only taken all together that these various borderline cases can reveal what it is like to experience one's body as one's own. Another of your arguments from phenomenal contrast in favor of the existence of a phenomenology of bodily ownership relies on evidence provided by the so-called "rubber hand illusion". The rubber hand illusion is a bodily illusion in which a participant sits in front of a fake hand aligned with one's body, while their real hand is hidden behind a screen (Botvinick \& Cohen, 1998). Their real hand is subsequently stroked while the fake hand is stroked either synchronously or asynchronously. Thus, participants feel tactile sensations on their real hand while they see the fake hand being stroked at the same time or with some delay. After a few minutes, most participants report the following effects in the synchronous condition but not in the asynchronous condition: (a) the tactile sensation of being stroked feels as if it was located on the rubber hand in front of them rather than on their real hand; and (b) it feels as if the rubber hand was their own hand. These reports typically come in the form of ratings of questionnaire items such as "It seems to me as if the rubber hand were my own hand". You have argued that bodily experiences elicited by the experimental setup of the rubber hand illusion in the synchronous condition and in the asynchronous condition, respectively, exhibit a phenomenal contrast that lends support to the view that there is a phenomenology of bodily ownership over the fake hand in the former condition. Questionnaire ratings are the primary source of evidence regarding this phenomenal contrast. As you acknowledge yourself, however, the difference in ratings of items related to ownership of the fake hand between asynchronous and synchronous conditions is generally not very impressive: on average, participants are barely in agreement with the idea that the fake hand seemed as if it was theirs during the synchronous condition (Vignemont, 2018a, p. 17; see Longo et al., 2008). Furthermore, recent work by Peter Lush and colleagues suggests that the rubber hand illusion does not adequately control for demand characteristics - the cues that may convey to participants the experimental outcome or response that the experimenter expects or desires (Lush et al., 2019; Lush, 2020; Roseboom \& Lush, 2020) Demand characteristics can not only influence participants' behavior, but also change their experience. Thus, expectancies arising from demand characteristics might cause participants to exert - unknowingly - top-down control of phenomenology, similarly to how they might respond to imaginative suggestions. In fact, Lush and colleagues Vignemont - Body Meet Self ALIUS Bulletin n4 (2020) aliusresearch.org/bulletin 118 found a substantial relationship between trait hypnotisability and both implicit (behavioral) and explicit (verbal) measures of the rubber hand illusion (Lush et al., 2019). In follow-up research, they found that participants' expectancies for "control" and "illusion" statements in synchronous and asynchronous conditions of the rubber hand illusion differ similarly to published illusion reports, implying that standard rubber hand illusion control measures do not effectively control for demand characteristics (Lush 2020). These findings suggest that ownership ratings in the rubber hand illusion may reflect implicit imaginative suggestion effects (in line with Alsmith, 2015). What do you make of these observations about the rubber hand illusion? Do you think they weaken the evidential strength of reports from the rubber hand illusion for your argument? When the rubber hand illusion started to be systematically tested, the community thought that we had the magic key to experimentally investigate the sense of body ownership. Since then, after more than 20 years of research and hundreds of versions, I believe that we are all less enthusiastic and we wish we could find other experimental paradigms to test bodily self- awareness. Still, despite their limits and weaknesses, it does not mean that we should just give up on bodily illusions. There are lessons to draw from all the results. We just have to be cautious. It is also interesting to note that there is a strong parallel with the literature on agency, we are just a few years late. As for agency, the influence of top- down factors is more important than we originally thought. With one of my students, Clement Apelian, we have actually compared the respective impact of hypnosis and sensory manipulation for bodily awareness (Apelian \& Vignemont, in preparation). Now the fact that an illusion can be influenced by cognitive factors does not entail that the illusion is constitutively cognitive. It just shows that it can be cognitively penetrated, possibly through imagination. Again, it is important to be vigilant about the significance of empirical results. There is a strong parallel with the literature on agency, we are just a few years late. " " Vignemont - Body Meet Self ALIUS Bulletin n4 (2020) aliusresearch.org/bulletin 119 According to your affective account of the phenomenology of ownership, the phenomenology of bodily ownership is a type of affective phenomenology that bears some resemblance with the feeling of familiarity one may have upon seeing well-known faces, such as those of family members and friends. Nonetheless, the affective feeling of bodily ownership is meant to be more specific than the generic feeling of familiarity. Indeed, unlike the latter, it exclusively tracks one object (one's own body), and it has a positive valence that is motivating for action (i.e., it motivates oneself to protect one's body) (Vignemont, 2018a, p. 192). For this reason, you describe this affective phenomenology as a "narcissistic" feeling, namely an awareness of the special significance that one's body has for oneself. As you also put it, in having a bodily experience "one is aware of bodily boundaries as having a special significance for the self" (Vignemont, forthcoming). This affective account of the phenomenology of bodily ownership is intended to address a dilemma raised by Peacocke (2015), that you summarize as follows: "if the protective body map represents one's body qua one's own, then it presupposes what it is supposed to explain, but if it does not, then one is left with no explanation of the first-personal character of the sense of bodily ownership" (Vignemont, 2018a, p. 204). To escape the first horn of this dilemma, you argue that bodily experiences that involve a phenomenology of bodily ownership do not explicitly represent the subject's body as their own. Instead, the "narcissistic" feeling of ownership represents the subject's body as the body that matters (or the body that has special significance) (Vignemont, forthcoming). Presumably, the "first-personal character" of the phenomenology of bodily ownership cannot refer to it having de se content, since it does not represent the subject's body as their own. Accordingly, you argue that the first-personal character of the phenomenology of bodily ownership is guaranteed by the "format" or "structure" of narcissistic feelings, rather than by their content. Indeed, you claim that the brain has a "protective body map" that always and exclusively tracks one's own body; as a result, bodily experiences anchored in the protective body map are automatically tagged with a sort of "self-centred glow" (Vignemont, 2018a, p. 205) derived from the body-tracking function of the protective body map. On your view, this "self-centred glow" is not a matter of the bodily experience having de se (e.g., first-personal) content, but is rather like the perspectival structure of visual experience: when one sees a tree in Vignemont - Body Meet Self ALIUS Bulletin n4 (2020) aliusresearch.org/bulletin 120 front of oneself, you suggest that one's visual experience represent the tree as being "in front" (simpliciter), rather than "in front of me". Similarly, when I feel pressure on my hand, my experience represents pressure "on the body that matters", rather than pressure "on my own body" (ibid., p. 205). One might wonder to what extent this analogy illuminates the nature of the "first-personal character" of the phenomenology of bodily ownership. It is certainly plausible that the sensory experience of very simple organisms, and perhaps even some visual experiences in humans, have merely de hinc spatial content without de se content - e.g., representing something as being in front of here, where here refers to the location of the point of origin of the sensory experience's egocentric frame of reference (see Peacocke, 2014, 2017; Schellenberg, 2016). However, many authors agree that ordinary visual experiences represent the locations of environmental landmarks with respect to the subject's own location as such; in other words, when one sees a tree, one's visual experience has a content of the type <There is a tree in front of me> (e.g., Cassam, 1997; Noe, 2005; Peacocke, 1998, 2014; Schwenkler, 2014). On this view, visual experiences typically have (nonconceptual) de se content. If that wasn't the case, then it might be difficult to see in what sense visual experiences have "first-personal character" at all - any more than a self- driving car's representation of its environment from the data provided by its sensors (cameras or LIDAR system) has "first-personal character", simply because it represents spatial properties of the environment in an egocentric frame of reference. Similarly, if bodily experiences represent one's body as the body that matters simpliciter, one might wonder where the "first-personal character" of such experiences comes from. To borrow an example from Martin (1995), what would be the difference between having a bodily experience with the content <There is hurt in this ankle>, which surely has no "first-personal character", and a bodily experience with the content <There is hurt in the ankle that matters>? On a related note, mattering and having a special significance are normally dyadic predicates: something can only matter or have a special significance to or for someone. Consequently, one might also wonder whether an experience can represent one's body as the body that matters without representing it as the body that matters to me (or, indeed, to someone else). What do you make of these concerns regarding the idea that the phenomenology of bodily ownership has "first-personal character" without de se content? Vignemont - Body Meet Self ALIUS Bulletin n4 (2020) aliusresearch.org/bulletin 121 To answer Peacocke's dilemma, which is about the body representation, all I need to claim is that the protective body map has no de se content. That's not the same as to say that the experience of ownership has no de se content. I actually do not make this latter claim. What I claim is that the relation to the subject should be understood in terms of personal significance, and not of myness. Actually, in my latest paper on ownership in the Journal of the APA (Vignemont, 2020), I argue that ownership experiences do have de se content. More specifically, I defend the view that the subject is part of the truth conditions of the content. Now one can defend this view and still leave implicit the self component. Such a move has been made by Perry (1993) about egocentric experiences, for instance. Egocentric terms actually illustrate how dyadic predicates can be represented as being monadic (Campbell, 1994). However, there is an alternative worth exploring, namely, that the de se nature of the experience of ownership follows from the format of bodily experiences instead of its content. This can take two forms. The first option is that there is a de se mode of bodily experiences. This is in line with Recanati's (2012) theory of mental files. The second option is that bodily experiences that can ground ownership judgments have a distinctive affective mental paint. This is in line with the attitude-based approach to affective representations (Mitchell, 2019; Deonna \& Teroni, 2012). To conclude, the main objective of the Bodyguard hypothesis was to develop the idea that the phenomenology of ownership is about personal significance. But a lot remains to be done to explain in detail how to analyse this notion of personal significance, at the level of content or of attitude. In a recent post on The Brains Blog (Vignemont, 2019b), you further elucidate the first horn of the dilemma raised by Peacocke (2015) as follows: "One may claim that the conceptual mineness content instantiated by the ownership judgment is grounded in a non-conceptual mineness content at the level of the feeling. But this solution seems to simply beg the question and to leave us with no explanation of the mineness content (conceptual or non-conceptual)." Vignemont - Body Meet Self ALIUS Bulletin n4 (2020) aliusresearch.org/bulletin 122 To understand the pull of this objection, it might be helpful to say something more about the nature of the required explanation, namely - what it is about the existence of a (non-conceptual) representation of a body part as one's own that begs for an explanation? Is what we are after an explanation of the origins of such representation through evolution and individual development, or a mechanistic explanation of why such representations feature in the content of any specific bodily experience, or both? Might the kind of explanation you provide for the existence of a representation of a body part as having a special significance not account equally well for the existence of a representation of a body part as one's own? The key question is what the difference is between a creature that can represent only "this leg is bent" (Peacocke's degree zero of self-representation) and a creature that nonconceptually represents "my leg is bent" (degree one of self-representation). To reply that the difference is only that at degree one the creature has a non-conceptual grasp of myness does not bring us very far. It is possible to posit myness as an irreducible primitive phenomenal property but I believe we should do so only when all the other attempts have failed. It seems more interesting to try to understand what new abilities are available to the creatures at degree one independently of myness. Another way to ask the question is what is required for one to be aware of one's body as one's own. I have argued that there are at least three abilities: (i) the spatial ability to individuate the boundary of the body; (ii) the affective ability to ascribe a value to one side of the body boundaries; and (iii) a general self- referential ability. It might then seem that I am also begging the question since I appeal to self-referential abilities in the explanans. However, this would not be true. My recipe is <personal significance + self = ownership>. I do not pretend to offer an explanation of self-awareness in general. My objective is more modest: I only offer an account of bodily self-awareness, this requires the creature to already have the ability to be self-aware. There is thus no circularity. A prediction of my view is thus that an individual who suffers from a complete disruption of self-referential abilities in other domains would not be able to experience her body as her own. This might be the case in depersonalization disorder, in which their sense of body disownership may result from a more general disruption. I do not pretend to offer an explanation of self- awareness in general. My objective is more modest: I only offer an account of bodily self- awareness. " " Vignemont - Body Meet Self ALIUS Bulletin n4 (2020) aliusresearch.org/bulletin 123 In the debate over the existence of a phenomenology of bodily ownership, at least four pairs of labels have been used to distinguish between available positions: (a) "realism" is opposed to "eliminativism" (Gallagher, 2017), (b) "inflationism" is opposed to "deflationism" (Bermudez, 2011; Vignemont, 2013; Gallagher, 2017, 2019; Serrahima, 2019); (c) "antireductionism" is opposed to "reductionism" (Martin, 1995); and (d) "liberalism" is opposed to "conservatism" (Vignemont, 2018a). In your own work, you have labelled your own view as "liberal" (Vignemont 2018a, p. 13) and "deflationary" (Vignemont, forthcoming). In so far as you defend the existence of a phenomenology of bodily ownership, your position can be adequately described as "realist"; and in so far as you argue that the phenomenology of ownership is not "an irreducible mineness quality" (Vignemont 2018a, p. 48), but is instead reducible to an affective quality, your view can also be presumably described as "reductionist". How do you understand these various labels, as qualifications of your view of the phenomenology of bodily ownership? Relatedly, do you think that some or all of the four dichotomies listed above are equivalent, or partially overlap, or are orthogonal? I think that Bermudez's original distinction between deflationary and inflationary views played an important role by starting a debate at the beginning. However, I've struggled for many years to exactly understand what he meant by it. I originally thought that inflationists claim that there is a phenomenology of ownership while deflationists claim that there is no such thing. But this is not the issue and both sides, as long as they are not eliminativist, agree that it feels like something to be aware of one's body as one's own. Nor was the issue whether one could give a reductionist account of the phenomenology of ownership. Indeed, since the beginning Bermudez acknowledged that my theory was reductionist and yet I was called an inflationist. The crucial question then is whether one defends the existence of a feeling of myness or not. Indeed, all the objections that Bermudez offers against inflationism are about myness, and only about that. In brief, inflationists are pro-myness and deflationists are anti-myness. However, in my first paper on ownership (Vignemont, 2007), I did not even use the term of myness (I double-checked), though it was on the basis of this paper that Bermudez decided that I was an inflationist. More explicitly, in Mind the body (2018a), I clearly criticize the myness hypothesis. On Bermudez's Vignemont - Body Meet Self ALIUS Bulletin n4 (2020) aliusresearch.org/bulletin 124 taxonomy, I am thus a deflationist. But most people are actually deflationist and I am not convinced that this taxonomy really helps understanding the disagreements between the various views. I find the classic distinction between liberal and conservative conceptions more informative. It brings the debate on ownership within the wider issue of the admissible content of perception. The crucial question is whether one defends the idea that bodily experiences represent only low-level somatosensory properties (pressure, posture, location, temperature, etc.) or whether they can represent high-level properties. These other properties include myness but not only. It can also include agentive properties and affective properties, such as personal significance and value. In brief, liberals are pro rich bodily content whereas conservatives are anti rich bodily content. Within this taxonomy, I am a liberal. So, to conclude, I am a liberal deflationist. In your view, pain and touch play complementary roles in the development of the sense of bodily ownership in healthy individuals (Vignemont, 2016, 2017b, 2018b). In a nutshell, touch makes an essentially spatial contribution, whereas pain makes an essentially affective contribution. In touch, we perceive non- bodily objects by being in contact with them, which makes the bodily boundaries especially salient. In virtue of this, touch fundamentally contributes to the individuation of the body with respect to other objects - namely those that fall outside the bodily boundary. In turn, pain adds an affective valence to the body that stands at one of the sides of the perceived boundary, contributing a phenomenology of bodily ownership to bodily sensations best described in terms of care or import, as discussed above. Roughly put, it is only the body that hurts. In your own words, pain "vividly highlights for the subject that what is inside bodily boundaries matters for the self, for its needs, its comfort, and its preservation" (Vignemont, 2017b, p. 475). Protective behavior is then one central manifestation of the sense of bodily ownership. For a start, is this a fair summary of your ideas on how touch and pain interact to give raise to a sense of bodily ownership? This is a perfect summary. The skin is a natural boundary that one can be aware of through touch. However, bodily self-awareness cannot be reduced to the spatial awareness of the body. We need more than the boundaries that touch can provide. We need to know which side of the boundaries we are in. And that's where pain, with its affective component, can play a role. Vignemont - Body Meet Self ALIUS Bulletin n4 (2020) aliusresearch.org/bulletin 125 Your description of how touch contributes to bodily ownership follows Mike Martin's (1992, 1993, 1995). Martin argued that all located bodily sensations convey a sense of the boundaries of the body. Arguably, his own analysis of all bodily sensations as involving a sense of what is inside and what is outside the bodily boundaries is an extension, to somatosensation in general, of the model that he endorses for tactile perception. In a paper with Olivier Massin (Vignemont \& Massin, 2015), you have called Martin's model of tactile perception the template model of tactile perception. However, you argue against the template model, and in favour of an alternative body map model of tactile perception. Does your account of bodily ownership rely in any way on the specificity of the body map model of touch, vis a vis the template model? Or do you see these two parts of your research as relatively unrelated? I am greatly indebted to Martin's work on touch and on the sense of ownership, though we do not reach the same conclusions. The template theory and the body map theory actually do not address the same dimensions of touch, the template theory focusing on the exteroceptive content of touch (e.g. how I feel the circular shape of the glass), and the body map theory focusing on its bodily content (e.g. how I feel pressure on my skin). For the sense of body ownership, what is directly relevant is the fact that the property of pressure is relational: it involves a force exerting on your body, which is independent of you. In a new work with Olivier Massin, we argue that touch gives us a unique sense of reality of what is touched because it presents it as being both mind-independent and causally efficient (Massin \& Vignemont, 2020). Touch informs us that the felt object can move, or have an effect, on other entities. When one actively touches an object, one exerts a force on it and one feels not only one's effort but also its resistance to one's effort: "There is no commoner remark than this, that resistance to our muscular effort is the only sense which makes us aware of a reality independent from ourselves" (James, 1890). Only effortful touch presents us Only effortful touch presents us with the contrast between ourselves as striving agents and an independent causally empowered being that resists our effort. Tactile experiences thus give us the boundary between body and world. " " Vignemont - Body Meet Self ALIUS Bulletin n4 (2020) aliusresearch.org/bulletin 126 with the contrast between ourselves as striving agents and an independent causally empowered being that resists our effort. Tactile experiences thus give us the boundary between body and world, which is required to draw the boundary of the body, a lot more than the visual experiences of the body because only tactile experiences can give this sense of reality of what is not the body. In your paper "The first person in pain" (Vignemont, 2018b), you write the following: "I would like to suggest that provided a pain is felt as one's own, one will react to it normally, even if the body part in which it is felt is not itself felt as one's own. Feeling pain is always of great concern to me, no matter where I feel it." In this quote, you leave room for the idea that normal reactions to pain, which include protective behaviour, result from the fact that pain be felt as one's own mental state, independently of whether or not one has a sense of ownership for the painful body. This idea, however, seems somehow in tension with your explanation of the development of the sense of bodily ownership. Protective behavior, and the affectivity attached to it, seem constitutive of what it is to experience the body as one's own, in the framework of your view. If the very subjectivity of pain experiences - the fact that they are felt as the subject's experiences - suffices to motivate normal protective behavior, does this not undermine the importance of the inside- outside distinction, provided by touch, for the sense of ownership? Imagine having a terrible pain in the back. You feel the urge to take a painkiller. This can be conceived as a protective behavior. However, the question is what you protect: your body or your psychological life? Let's imagine now that your pain was actually indicative of a kidney infection. By hiding the symptom, you were actually delaying the diagnosis and thus, the cure of the disease. Many have discussed this killing-the-messenger problem and asked whether it is rational or not to take painkillers. But my point here is simply that there are protective behaviours that primarily concern our mental life, and not our body, and that pain, like other negative emotions, can also elicit this type of response, no matter where we locate our pain. I do not think that we should take protective behavior as being a unified category. For instance, one may say that to put money on the side to make sure you Vignemont - Body Meet Self ALIUS Bulletin n4 (2020) aliusresearch.org/bulletin 127 still have something to survive on when you retire is a protective behavior. But it clearly does not involve the protective body map. As noted in the last chapter of Mind The Body, one should be extremely careful when analyzing the way people protect or fail to protect themselves. Even in the case of pain, the relation to pain responses is complex and deserves a detailed treatment. One topic you have worked extensively on is the nature and cognitive underpinnings of vicarious pain. In a very rich paper with Pierre Jacob (Vignemont \& Jacob, 2012), you define different ways in which we may feel the pains of other people. You defend that vicarious pain consists in imagining being in pain - in particular, doing so through Enactment-imagination (E- imagination): we mentally simulate the psychological state of the other by activating our pain system offline. On your proposal, there are two variants of vicarious pain, depending on which of the two subsystems actually involved in pain processing are activated when E-imagining pain. On the one hand, the (offline) activation of the sensory-discriminative subsystem gives rise to a form of pain contagion: a self- centered type of vicarious pain, in which one imagines how it would be to feel pain in the bodily location of one's own body that maps the other's, and anticipates the sensorimotor consequences of such pain. On the other hand, the (offline) activation of the affective subsystem gives rise to empathetic pain: a global, non-localised bodily feeling, directed to the other's body. In sum, when it comes to vicarious pains, we have (i) pain-like bodily feelings that presumably one identifies as one's own feelings; (ii) caused by pain that affects the bodies of others; (iii) at least in the case of empathetic pain, actually directed to the bodies of others (that is, not clearly centered on what it is like, or what it would be like, to feel a certain pain on one's own body); (iv) yet, and again especially in the case of empathetic pain, involving affective reactions directed to the bodies of others. In your view, how do these features of vicarious pain interact with the role you ascribe to pain in the emergence of the feeling that a body is one's own? Do we have, for the bodies of others, any feeling similar in any relevant respect to those constitutive of the sense of bodily ownership? If so, how do you reconcile this with the fact that, in normal conditions, the sense of bodily ownership tracks only one's own body? One of the take home messages of our work with Pierre Jacob is that vicarious pain is not the same as standard pain. No matter how much we might say 'I Vignemont - Body Meet Self ALIUS Bulletin n4 (2020) aliusresearch.org/bulletin 128 feel your pain', the fact is that we don't. Part of the difference comes from the content, as you rightly noted. Part of it comes from the fact that it is under the imaginative mode. Thanks to these differences, one does not normally confuse other people's pain with one's own pain. As a consequence, one does not react to one's vicarious pains in the same way as to one's standard pain. Vicarious pains have thus little consequences for the sense of body ownership. Contagious pain remains self-centered, all about one's body. Empathetic pain, on the other hand, is other-centered but it is almost disembodied. In a recent talk at the Ecole Normale Superieure, entitled "Keeping the world at distance", you addressed the notion of social distancing, which has become central to our lives in current pandemic times. In the talk you discussed how social distancing reveals the sensorimotor mechanism of peripersonal space, namely the space immediately surrounding our body. You have argued in print that, because anything occupying peripersonal space might soon be in touch with the body, peripersonal perception has evolved into a type of perception distinctively linked to protective action (Vignemont, 2018c). The notion of peripersonal space thus meshes nicely with your ideas on bodily ownership as bodily care: in a way, the space around our bodies is not completely alien, and our sensorimotor system must be sensitive to it just as it is sensitive to events on or within the bodily boundaries. However, social distancing is not all about keeping the right distance from others in order to promote self-preservation. We also keep apart from the others in order to protect them from ourselves. Do you think that your notions of peripersonal space and self-care have any implications in this direction, now that attention to other-preservation determines our social interactions? Most of my recent research has indeed been on peripersonal space and we have just edited the first multidisciplinary volume on the topic with Alessandro Farne, Andrea Serino and Hong-Yu Wong, entitled The World At Our Fingertips (Vignemont et al. 2021). Peripersonal space is a fascinating notion that is the crossroad between perception and action and between self and others. It has an immediate relevance with what is happening right now in the world. Now it is a real open question whether social distancing is more about self-preservation or other-preservation. One can easily imagine that altruism increases social distancing, but could it suffice for social distancing? Imagine you could not catch Covid-19 but you could still contaminate others. How likely is it that your distancing behavior will be more reflexive and less Vignemont - Body Meet Self ALIUS Bulletin n4 (2020) aliusresearch.org/bulletin 129 spontaneous than if you could also get contaminated yourself? But maybe I should trust more in human nature. What do you think will be hot topics in the philosophical and empirical literature on bodily awareness in the coming years? More and more people have been talking about virtual reality. However, this is just a tool, a very convenient one, but without the right questions, it won't bring us very far. Many also work on interoception. But though there is something quite fascinating about the constant flow of inner signals we receive, we still need to first provide a theory of interoception before trying to assess what impact it has for our cognitive life. To tell me that high interoceptive score is correlated with this or that does not tell me what interoception is. Is interoception a natural kind or just an umbrella term to refer to many different inner signals? What does interoceptive awareness consist in? Is it perceptual or not? Can it ground knowledge? Does it have a distinctive phenomenology? Can one give a representationalist account of interoceptive experiences? What properties would they then represent? As far as I know, no philosopher has yet tried to answer these questions. Another fascinating topic is augmentative technology. Sci-fi movies show our future selves endowed with high-tech prostheses and exoskeleton. But could we actually exploit a third artificial arm at no cost for our biological arms? Does the brain process artificial devices more like tools or more like body parts? Or could they have their own sensorimotor and phenomenological signature? And how to best design these artificial devices? These are some questions for the future and I'm sure new ones will arise. References Alsmith, A. (2015). Mental Activity and the Sense of Ownership. Review of Philosophy and Psychology, 6(4), 881-896. https://doi.org/10.1007/s13164014-0208-1 Apelian, C., \& de Vignemont, F. (in preparation). Hypnotic distortions of body representations. Vignemont - Body Meet Self ALIUS Bulletin n4 (2020) aliusresearch.org/bulletin 130 Bayne, T. (2009). Perception and the Reach of Phenomenal Content. The Philosophical Quarterly, 59(236), 385-404. https://doi.org/10.1111/j.1467-9213.2009.631.x Bermudez, J. L. (2011). Bodily Awareness and Self-Consciousness. In S. Gallagher (Ed.), The Oxford Handbook of the Self. Oxford University Press. Billon, A., \& Kriegel, U. (2015). Jaspers' Dilemma: The Psychopathological Challenge to Subjectivity Theories of Consciousness. In R. J. Gennaro (Ed.), Disturbed Consciousness: New Essays on Psychopathology and Theories of Consciousness (pp. 29-54). MIT Press. Bottini, G., Bisiach, E., Sterzi, R., \& Vallarc, G. (2002). Feeling touches in someone else's hand. NeuroReport, 13(2), 249. Botvinick, M., \& Cohen, J. (1998). Rubber hands 'feel' touch that eyes see. Nature, 391(6669), 756-756. https://doi.org/10.1038/35784 Campbell, J. (1994). Past, Space, and Self. MIT Press. Cassam, Q. (1997). Self and World. Oxford University Press. Chadha, M. (2018). No-Self and the Phenomenology of Ownership. Australasian Journal of Philosophy, 96(1), 14-27. https://doi.org/10.1080/00048402.2017.1307236 Chudnoff, E. (2015). Phenomenal Contrast Arguments for Cognitive Phenomenology. Philosophy and Phenomenological Research, 91(1), 82-104. https://doi.org/10.1111/phpr.12177 Deonna, J., \& Teroni, F. (2012). The Emotions: A Philosophical Introduction. Routledge. Evans, G. (1982). The Varieties of Reference. Oxford: Oxford University Press. Folegatti, A., Farne, A., Salemme, R., \& de Vignemont, F. (2012). The Rubber Hand Illusion: Two's a company, but three's a crowd. Consciousness and Cognition, 21(2), 799-812. https://doi.org/10.1016/j.concog.2012.02.008 Gallagher, S. (2017). Deflationary accounts of the sense of ownership. The Subject's Matter: Self-Consciousness and the Body. The MIT Press, Cambridge MA, 145-162. Gallagher, S. (2019). The Senses of a Bodily Self. ProtoSociology, 36, 414-433. https://doi.org/10.5840/protosociology20193616 Gouzien, A., de Vignemont, F., Touillet, A., Martinet, N., De Graaf, J., Jarrasse, N., \& Roby-Brami, A. (2017). Reachability and the sense of embodiment in amputees using prostheses. Scientific Reports, 7(1), 4999. https://doi.org/10.1038/s41598-017-05094-6 James, W. (1890). The principles of psychology. Henry Holt and Company. http://content.apa.org/books/11059-000 Vignemont - Body Meet Self ALIUS Bulletin n4 (2020) aliusresearch.org/bulletin 131 Longo, M. R., Schuur, F., Kammers, M. P. M., Tsakiris, M., \& Haggard, P. (2008). What is embodiment? A psychometric approach. Cognition, 107(3), 978- 998. https://doi.org/10.1016/j.cognition.2007.12.004 Lush, P. (2020). Demand Characteristics Confound the Rubber Hand Illusion. Collabra: Psychology, 6(1), 22. https://doi.org/10.1525/collabra.325 Lush, P., Botan, V., Scott, R. B., Seth, A., Ward, J., \& Dienes, Z. (2019). Phenomenological control: Response to imaginative suggestion predicts measures of mirror touch synaesthesia, vicarious pain and the rubber hand illusion. https://doi.org/10.31234/osf.io/82jav Makin, T. R., de Vignemont, F., \& Micera, S. (2020). Soft Embodiment for Engineering Artificial Limbs. Trends in Cognitive Sciences, 24(12), 965-968. https://doi.org/10.1016/j.tics.2020.09.008 Martin, M. G. F. (1992). Sight and Touch. In T. Crane (Ed.), The Contents of Experience (pp. 199-201). Cambridge University Press. Martin, M. G. F. (1993). Sense modalities and spatial properties. In N. Eilan, R. McCarty, \& B. Brewer (Eds.), Spatial representation (pp. 206-218). Oxford University Press. Martin, M. G. F. (1995). Bodily Awareness: A Sense of Ownership. In J. L. Bermudez, A. J. Marcel, \& N. M. Eilan (Eds.), The Body and the Self (pp. 267-289). MIT Press. Massin, O., \& de Vignemont, F. (2020). Unless I Put My Hand into His Side, I Will Not Believe. In D. E. Gatzia \& B. Brogaard (Eds.), The Epistemology of Non- Visual Perception (p. 165). Oxford University Press. Mitchell, J. (2019). Affective Representation and Affective Attitudes. Synthese, 1- 28. https://doi.org/10.1007/s11229-019-02294-7 Noe, A. (2005). Action in Perception. MIT Press. Peacocke, C. (1998). Being Known. Oxford University Press. Peacocke, C. (2014). The Mirror of the World: Subjects, Consciousness, and Self- Consciousness. Oxford University Press. Peacocke, C. (2015). Perception and the First Person. In M. Matthen (Ed.), The Oxford Handbook of Philosophy of Perception. http://www.oxfordhandbooks.com/view/10.1093/oxfordhb/9780199600472. 001.0001/oxfordhb-9780199600472-e-022 Peacocke, C. (2017). Philosophical Reflections on the First Person, the Body, and Agency. In F. de Vignemont \& A. Alsmith (Eds.), The Subject's Matter: Self- Consciousness and the Body (pp. 289-310). MIT Press. Perry, J. (1993). The Problem of the Essential Indexical: And Other Essays. Oxford University Press. Vignemont - Body Meet Self ALIUS Bulletin n4 (2020) aliusresearch.org/bulletin 132 Recanati, F. (2012). Mental Files. Oxford University Press. Romano, D., \& Maravita, A. (2019). The dynamic nature of the sense of ownership after brain injury. Clues from asomatognosia and somatoparaphrenia. Neuropsychologia, 132, 107119. https://doi.org/10.1016/j.neuropsychologia.2019.107119 Roseboom, W., \& Lush, P. (2020). Serious problems with interpreting rubber hand illusion experiments. https://doi.org/10.31234/osf.io/uhdzs Schellenberg, S. (2016). De Se Content and De Hinc Content. Analysis, 76(3), 334- 345. https://doi.org/10.1093/analys/anw019 Schwenkler, J. (2014). Vision, Self-Location, and the Phenomenology of the 'Point of View.' Nous, 48(1), 137-155. https://doi.org/10.1111/j.1468-0068.2012.00871.x Serrahima, C. (2019). My Body is the Subject's Body. In Defence of Experientialism about the Sense of Bodily Ownership [University of Barcelona]. http://diposit.ub.edu/dspace/handle/2445/137017 Siegel, S. (2007). How Can We Discover the Contents of Experience? The Southern Journal of Philosophy, 45(S1), 127-142. https://doi.org/10.1111/j.2041-6962.2007.tb00118.x Siegel, S. (2010). The Contents of Visual Experience. Oxford University Press USA. Strawson, G. (1994). Mental Reality. MIT Press. Vallar, G., \& Ronchi, R. (2009). Somatoparaphrenia: A body delusion. A review of the neuropsychological literature. Experimental Brain Research, 192(3), 533- 551. https://doi.org/10.1007/s00221-008-1562-y de Vignemont, F. (forthcoming). The phenomenology of bodily ownership. In M. Garcia-Carpintero \& M. Guillot (Eds.), The sense of mineness. Oxford University Press. de Vignemont, F. (2007). Habeas Corpus: The Sense of Ownership of One's Own Body. Mind \& Language, 22(4), 427-449. https://doi.org/10.1111/j.1468-0017.2007.00315.x de Vignemont, F. (2011). Bodily immunity to error. In S. Prosser \& F. Recanati (Eds.), Immunity to error through misidentification: New essays (pp. 1-27). Cambridge University Press. de Vignemont, F. (2013). The mark of bodily ownership. Analysis, 73(4), 643-651. https://doi.org/10.1093/analys/ant080 de Vignemont, F. (2016). Pain and the Spatial Boundaries of the Bodily Self. In L. Garcia-Larrea (Ed.), Pain and the Conscious Brain. IASP. Vignemont - Body Meet Self ALIUS Bulletin n4 (2020) aliusresearch.org/bulletin 133 de Vignemont, F. (2017a). Can I see your pain? An evaluative model of pain perception. In J. Corns (Ed.), Can I see your pain? An evaluative model of pain perception (pp. 255-265). Routledge. de Vignemont, F. (2017b). Pain and Touch. The Monist, 100(4), 465-477. https://doi.org/10.1093/monist/onx022 de Vignemont, F. (2018a). Mind the Body: An Exploration of Bodily Self- Awareness. Oxford University Press. de Vignemont, F. (2018b). The First-Person in Pain. In M. Brady, D. Bain, \& J. Corns (Eds.), Philosophy of Pain. Unpleasantness, Emotion, and Deviance. Routledge. de Vignemont, F. (2018c). Peripersonal perception in action. Synthese. https://doi.org/10.1007/s11229-018-01962-4 de Vignemont, F. (2019a). Against Phenomenal Parsimony: A Plea for Bodily Feelings. In A. I. Goldman \& B. P. McLaughlin (Eds.), Metaphysics and Cognitive Science (pp. 268-283). Oxford University Press. de Vignemont, F. (2019b). Mind the body (4) What kind of first-personal content? In The Brains Blog. http://philosophyofbrains.com/2019/02/28/mind-the- body-4-what-kind-of-first-personal-content.aspx de Vignemont, F. (2020). What Phenomenal Contrast for Bodily Ownership? Journal of the American Philosophical Association, 6(1), 117-137. https://doi.org/10.1017/apa.2019.34 de Vignemont, F., \& Jacob, P. (2012). What Is It like to Feel Another's Pain? Philosophy of Science, 79(2), 295-316. https://doi.org/10.1086/664742 de Vignemont, F., \& Massin, O. (2015). Touch. In M. Matthen (Ed.), The Oxford Handbook of Philosophy of Perception (pp. 294-313). Oxford University Press. de Vignemont, F., Serino, A., Wong, H. Y., \& Farne, A. (forthcoming). Peripersonal space: A special way of representing space. In The world at our fingertips: A multidisciplinary investigation of peripersonal space. Oxford University Press. ISBN: 9780198851738\par
\endgroup

\ifdefined\ALIUSIssueBuild
\clearpage
\expandafter\endinput
\fi
\end{document}
